% Generated draft: requires mathematical and source review.
\documentclass[11pt]{article}
\usepackage[margin=25mm]{geometry}
\usepackage{fontspec}
\setmainfont{Latin Modern Roman}
\usepackage{amsmath,amssymb,mathtools,unicode-math}
\setmathfont{Latin Modern Math}
\usepackage{enumitem,needspace,xurl,hyperref}
\setlength{\parindent}{0pt}
\setlength{\parskip}{5pt}
\setlength{\emergencystretch}{3em}
\newcommand{\N}{\mathbb{N}}
\newcommand{\Z}{\mathbb{Z}}
\newcommand{\Q}{\mathbb{Q}}
\newcommand{\R}{\mathbb{R}}
\newcommand{\C}{\mathbb{C}}
\newcommand{\F}{\mathbb{F}}
\newcommand{\A}{\mathbb{A}}
\newcommand{\PP}{\mathbb{P}}
\newcommand{\E}{\mathbb{E}}
\newcommand{\eps}{\varepsilon}
\newcommand{\dd}{\,\mathrm{d}}
\begin{document}
\section*{UnsolvedMath ID: 30000160}
\textbf{Source problem code:} catalogue entry
\subsection{Definitions and mathematical statement}
Put \(q=e^{2\pi i\tau}\), \(\operatorname{Im}\tau>0\). The source considers the normalized extremal modular form
\[
 f(\tau)=1+\sum_{i\geq1}a_iq^i
\]
of weight \(k=12\mu\). In the usual level-one interpretation, \(f\in M_{12\mu}(\operatorname{SL}_2(\Z))\) and extremality means the normalized maximal initial gap, \(a_1=\cdots=a_\mu=0\). The excerpt does not explicitly state the modular group or coefficient ring; integral, or at least \(p\)-integral, coefficients are needed for the congruences. These conventions must be confirmed rather than generalized to arbitrary levels.

The expanded ``Case (2)'' hypothesis for a prime \(p\) is
\[
 p\nmid i\Longrightarrow a_i\equiv0\pmod p\quad(i\geq1),
 \qquad\exists i\geq1:\ a_i\not\equiv0\pmod p.
\]
The second condition concerns positive Fourier coefficients, not the constant coefficient \(1\); it excludes a constant reduction. No restriction to odd primes or \(p\geq5\) is supplied.

Conjecture 4 has two retained clauses:
\begin{enumerate}
\item There exists a modular form \(g\), of smaller weight \(k'<12\mu\) in the intended modular-form setting, such that
\[
 f(\tau)\equiv g(p\tau)\pmod p.
\]
\item There exists an \emph{extremal} modular form \(g\) of smaller weight and a natural number \(r\) such that
\[
 f(\tau)\equiv g(p^r\tau)\pmod p.
\]
Here extremality is in \(g\)'s own weight space; that weight is not asserted to be a multiple of twelve. The nontrivial iterated-pullback reading takes \(r\geq1\); the record does not fix whether its natural-number convention allows zero. The forms in the two clauses need not be the same.
\end{enumerate}
Both congruences are coefficientwise at the Fourier expansion, with the relevant coefficients \(p\)-integral.

The hypothesis alone already gives a formal series \(h(q)=1+\sum_{j\geq1}\overline a_{pj}q^j\) over \(\F_p\) with \(\overline f(q)=h(q^p)\). The conjectural content is that a modular form of smaller weight can supply the series, and in the second clause can be chosen extremal after a suitable iterate. Formal support extraction by itself does not establish either modularity or weight decrease.

The supplied triage dated 21 August 2026 reports proofs in substantial cases and further restrictions, while emphasizing that the full assertion remains Conjecture 4 in the report. Those cases and hypotheses are not reproduced. The draft preserves that partial-resolution assessment without turning either clause into an unconditional theorem.
\subsection{Short English statement}
The coefficient condition forces the reduction to be a series in \(q^p\), but the conjecture asks for a lower-weight modular realization. Its second clause additionally requires extremality and allows an iterated pullback. Group, coefficient, and natural-number conventions require checking.
\subsection{Sources}
\begin{itemize}
\item Local UnsolvedMath registry, permanent ID \(30000160\), exact source code \texttt{catalogue entry}, Oberwolfach Reports Open Problems, Conjecture 4. Statement status \texttt{corrected\_verified}; supplied verification expands Case (2) into the coefficient condition and retains both clauses.
\item ``Gitter und Anwendungen (2005),'' Oberwolfach Reports, pp. 5--67, DOI \texttt{10.4171/owr/2005/01}: \url{https://doi.org/10.4171/owr/2005/01}. Proposed year 2005.
\item Supplied ``primary OWR paper text'': \url{https://oa.tib.eu/renate/bitstreams/b4dac21d-26f7-4914-a2d2-f12c31dc9e53/download}.
\item Supplied ``Koike, 2009'': \url{https://doi.org/10.2206/kyushujm.63.123}. The exact covered cases and full theorem statements are not supplied; none of these documents was fetched.
\item Literature checked 21 August 2026; supplied status \texttt{partially\_solved}. Registry created 25 August 2026 and updated 16 September 2026. This dated curation triage requires independent verification before citation.
\item Registry curation attribution: Ulam AI / UnsolvedMath Contributors, CC BY 4.0. Underlying report and literature retain their supplied provenance; this statement-only draft requires human review.
\end{itemize}
\end{document}
